\chapter{Parameterstudien und polynomiale Surrogate}\label{chap:surrogates}
\section{Abgegrenzter Daten- und Methodenvertrag}
M8-Studien erzeugen stationäre FEM-Fälle mit deterministischen IDs und
Model-/Netz-/Solver-Fingerprints. Wiederverwendung und Fortsetzung setzen
passende Herkunft und vollständige erfolgreiche Ergebnisse voraus.
Die Kennwerttabelle unterscheidet Einheiten, Rollen, Gültigkeit und
Bilanzstatus. Ein Surrogat approximiert daraus reelle skalare Ziele;
es löst bei Training oder Vorhersage selbst kein FEM-Modell.
Es gibt keine neuronalen Netze, Feldrekonstruktion, POD/ROM, inverse
Identifikation, Optimierung oder Unsicherheitsquantifizierung in dieser API.
\index{Parameterstudie}\index{Surrogat}

Der Datensatz besteht aus $\mathbf X\in\mathbb R^{n\times p}$,
$\mathbf Y\in\mathbb R^{n\times q}$, eindeutigen Fall-IDs und
Parameter-/Zielnamen mit Einheit und Definition. Eingaben sind endliche
reelle numerische Werte; Boolesche, komplexe oder fehlende Werte werden nicht
still kodiert oder imputiert. Bei Kennwerttabellen werden nicht erfolgreiche
Fälle, fehlgeschlagene Solves/Bilanzen und ungültige ausgewählte Werte
standardmäßig abgelehnt. Ein ausdrücklich erlaubter Ausschluss protokolliert
jede ID und alle Gründe. Nicht geprüfte Qualität gilt nicht als bestanden.

Trainings- und Test-IDs sind eine vollständige, disjunkte, nichtleere Partition.
Identische numerische Parametervektoren dürfen die Splitgrenze nicht
überschreiten. Duplikate in einer Gruppe zählen nicht als zusätzliche
unabhängige Stützstellen. Optionales Seed-Sampling permutiert eindeutige
Parametergruppen gemeinsam; gespeicherte endgültige IDs sind maßgeblich.
Grad und Ridge-Gewicht legt der Aufrufer vorher fest. Wiederholte Auswahl am
Testfehler verbraucht dessen Unabhängigkeit und verlangt neue Holdouts.

\section{Train-only-Skalierung und vollständige Basis}
Nur Trainingsdaten liefern Minima $l_j$ und Maxima $h_j$:
\begin{equation}\label{eq:surrogate-scale}
c_j=l_j/2+h_j/2,\qquad s_j=h_j/2-l_j/2,
\qquad z_j=(x_j-c_j)/s_j.
\end{equation}
Damit liegen Trainingswerte je Spalte in $[-1,1]$; $z$ ist dimensionslos.
Konstante oder numerisch nicht skalierbare Parameter sind Fehler und müssen
fachlich begründet entfernt werden. Konstante Zielwerte sind dagegen zulässig.
Center und Scale werden nicht aus Holdouts nachberechnet.

Die Grad-1-Basis ist $(1,z_1,\ldots,z_p)$, die Grad-2-Basis ergänzt alle
$z_i z_j$ mit $i\le j$ in lexikographischer Indexreihenfolge.
Die Anzahl Basisterme ist
\begin{equation}\label{eq:surrogate-basis}
m_1=1+p,\qquad m_2=1+p+p(p+1)/2.
\end{equation}
Die Designmatrix hat $A_{ra}=z_r^{\alpha_a}$ für gespeicherte Exponenten.
Mindestens $m$ verschiedene Trainingspunkte sind nötig, auch bei Ridge.
Diese Anzahl ist notwendig, aber nicht hinreichend: Ein geometrisch
kollineares Design kann trotzdem rangdefizient sein.

\section{Least Squares, Ridge und Rang}
Koeffizienten $\mathbf B\in\mathbb R^{m\times q}$ erfüllen
\begin{equation}\label{eq:ridge}
\min_{\mathbf B}\|\mathbf A\mathbf B-\mathbf Y_{\rm train}\|_F^2
+\lambda\|\mathbf B_{1:, :}\|_F^2,\qquad\lambda\ge0.
\end{equation}
Der konstante Term wird nicht bestraft. Mehrere Zielspalten werden unabhängig
gelöst; der numerische Frobeniusausdruck ist kein dimensionsübergreifender
physikalischer Gütewert. Koeffizienten je Ziel tragen dessen Einheit;
Parameter-Einheiten sind in der Skalierung enthalten.
Das Gewicht gilt für diese skalierte, nicht durch Fallzahl normalisierte
Zielfunktion. Eine Verdopplung von Trainingszeilen verändert daher bei gleichem
$\lambda$ die relative Regularisierungsstärke.
\index{Ridge}\index{Least Squares}

Mit $\mathbf R=\operatorname{diag}(0,1,\ldots,1)$ wird
\begin{equation}\label{eq:augmented-lstsq}
\begin{pmatrix}\mathbf A\\\sqrt\lambda\mathbf R\end{pmatrix}\mathbf B
\simeq\begin{pmatrix}\mathbf Y_{\rm train}\\0\end{pmatrix}
\end{equation}
per \texttt{numpy.linalg.lstsq} gelöst, ohne Normalengleichungen.
Die NumPy-Primärquelle \cite{numpy-lstsq} beschreibt Least-Squares-Lösung,
Singulärwerte und numerischen Rang. Im Projekt gilt deren Standard-Rangtoleranz;
es wird keine eigene versteckte Regularisierung gewählt.
Unregularisierter Rangmangel ist ein Fehler; Ridge kann kollineare Designs
stabilisieren, aber ein rangdefizientes erweitertes System wird ebenfalls
abgelehnt. Extremes Gewicht oder nicht darstellbare Polynome sind ungültig.
Eine Konditionsverbesserung ersetzt keine geeigneten Datenpunkte.

\section{Vorhersage, Testfehler und Grenzen}
Für eine neue Zeile berechnet man dieselbe Skalierung und Basis;
$\hat{\mathbf y}=\mathbf a(z)^T\mathbf B$ in der gespeicherten Zielreihenfolge.
Die Vorhersageausgabe hat immer Form $(n,q)$, auch bei einem einzigen Punkt.
Die Trainingsbox ist $[l,h]$ mit inklusiven Grenzen. Außerhalb wird standardmäßig
abgelehnt; ausdrücklich erlaubte Extrapolation bleibt zeilenweise markiert.
Es gibt kein Clipping, keine stillen Einheitenumrechnungen und keine
Toleranzerweiterung der Box.

Die Box ist nicht die konvexe Hülle der Trainingspunkte und keine
Genauigkeitsgarantie innerhalb des Bereichs. Datenlücken, Materialwechsel,
Solverfehler und schlecht approximierte Nichtlinearität bleiben Risiken.
Testpunkte außerhalb markieren den Bericht, vergrößern aber nicht die Box.
Empirische Fehler je Ziel werden getrennt für Train/Test berechnet:
\begin{equation}\label{eq:surrogate-errors}
{\rm MAE}=\frac1n\sum_r|e_r|,\quad
{\rm RMSE}=\sqrt{\frac1n\sum_r e_r^2},\quad
{\rm maxabs}=\max_r|e_r|,\quad e_r=\hat y_r-y_r.
\end{equation}
Alle haben die Zieleinheit. Es werden keine bei Nullzielen undefinierten
relativen Fehlerquotienten, Konfidenzintervalle oder garantierten
Fehlerschranken ausgegeben. Surrogatfehler gegen diskrete FEM-Kennwerte,
FEM-Diskretisierungsfehler und experimentelle Modellabweichung sind drei
verschiedene Fehlerquellen.

\section{Handrechnungen: lineare Skalierung und Ridge}
Trainingspunkte $x=(-2,-1,0,1,2)$ m und Ziele $y=3+2x$ K (Koeffizient
2 K/m) geben $c=0$, $s=2$ m, $z=x/2$ und $y=3+4z$.
Die Koeffizienten sind $(3,4)$ K. Unabhängige Testpunkte $x=0{,}5$ und
$1{,}5$ m geben exakt $4$ bzw. $6$ K. Ein zweiter Ausgang $y_2=7$ W
bleibt konstant, obwohl konstante Eingabeparameter verboten sind.
Der bestehende Test prüft diese Zahlen bis $10^{-12}$ absolut.

Für ein symmetrisches Zwei-Punkt-Design $z=(-1,1)$ und $y=(1,5)$,
$\lambda=2$ ist der Intercept $b_0=3$.
Aus $\partial_{b_1}[(b_0-b_1-1)^2+(b_0+b_1-5)^2+2b_1^2]=0$
folgt $b_1=1$ statt der unregularisierten Steigung 2.
Ein zusätzlicher Testpunkt $z=0$ bleibt bei 3; der Intercept ist nicht
gegen null geschrumpft. Das Beispiel erklärt die Wirkung eines ausdrücklich
gewählten Gewichts, ohne eine automatische Modellwahl zu behaupten.

Für Beispiel-Residuen $e=(1,-2)$ sind MAE $=1{,}5$,
RMSE $=\sqrt{5/2}$ und Maximalfehler $=2$ in der Zieleinheit.
Ein quadratisches Ziel wie $y=x^2$ kann durch eine unzureichende Grad-1-Basis
sichtbare Holdoutfehler erzeugen; ein kleiner Trainingsfehler beweist keine
allgemeine Vorhersagegüte.

\section{Persistenz und Herkunft}
Das JSON-Artefakt speichert Basis, Skalierung, Box, Koeffizienten,
Metadaten für Parameter und Ziele, Quellnachweis und Prüfbericht.
SHA-256 schützt gegen unbeabsichtigte Beschädigung, nicht absichtliche
Manipulation. Strukturell gültiges Laden rekonstruiert keine historischen
Trainingsnachweise. Die gespeicherten Train/Test-IDs, Seeds, Daten-/Tabellen-
Fingerprints und Versionen machen die numerischen Entscheidungen prüfbar.
Kein Pickle oder ausführbarer Modellcode wird deserialisiert.

\subsection*{Implementierung und unabhängige Nachweise.}
{\small\raggedright
\path{src/fem_toolbox/studies.py},
\path{src/fem_toolbox/study_reports.py},
\path{src/fem_toolbox/surrogates.py};
\path{tests/unit/test_surrogates.py},
\path{tests/regression/test_surrogate_example.py},
\path{examples/surrogate_heat/README.md}.
Die echten Wärme-FEM-Fälle im Beispiel werden zuerst erzeugt und gegen
unabhängige diskrete Erwartungen geprüft. Das Training selbst enthält
keinen verdeckten FEM-Lauf. Doc/59--61 bleiben die kanonischen API-Verträge.

\par}
